%% kapa: exam=1 round=36 year=2025 period=2 form=A subject=accounting subject_ko=회계학 seq=1 number=41
%% source: Q-Net 감정평가사 기출문제 (cst00309) artlSeq=5231525 file=raw/1008/5231525/2025년 제36회 감정평가사 1차시험 2교시 문제지 원본.pdf
%% conversion: kordoc 4.24.51 -> tools/gen1.py (auto); review=proofread-sonnet; flags=none
\begin{question}{41}
유용한 재무정보의 질적특성에 관한 설명으로 옳은 것은?
\begin{choices}
  \item 완벽한 표현충실성을 위해서 서술은 완전하고 중요하며 오류가 없어야 할 것이다.
  \item 재무정보가 예측가치를 갖기 위해서는 그 자체가 예측치 또는 예상치이어야 한다.
  \item 나타내고자 하는 바를 충실하게 표현하는 가장 목적적합한 정보를 선택하려는 결정의 결과가 비대칭성인 경우라도 특정 회계기준에서 비대칭적인 요구사항을 포함할 수 없다.
  \item 오류가 없다는 것은 현상의 기술에 오류나 누락이 없고 보고 정보를 생산하는 데 사용되는 측정과 절차 측면에서 완벽하게 정확하다는 것을 의미한다.
  \item 합리적인 추정치의 사용은 재무정보의 작성에 필수적인 부분이며 추정치가 명확하고 정확하게 기술되고 설명되는 한 정보의 유용성을 저해하지 않는다.
\end{choices}
\end{question}
\answer{5}
